% =====================================================================
% Chương 0 — Tổng quan hành trình dữ liệu (4 lớp IoT chuẩn ↔ repo)
% =====================================================================

\setcounter{chapter}{-1}

\chapter{Chương 0 --- Tổng quan hành trình dữ liệu}

Chương này đặt bối cảnh cho toàn bộ báo cáo: ánh xạ \textbf{4 lớp IoT chuẩn}
(Cảm biến, Mạng, Xử lý, Ứng dụng) vào hệ thống UMT-TBS-official và vẽ sơ đồ
dòng dữ liệu trên cả hai đường truyền. Các chương sau (1--7) đi sâu vào từng
thư mục, từng file, từng hàm và từng dòng code.

\section{Bốn lớp IoT chuẩn nhìn từ hệ thống}

\begin{center}
\small
\renewcommand{\arraystretch}{1.2}
\begin{longtable}{@{}p{3.4cm}p{4.6cm}p{4.2cm}p{4.2cm}@{}}
\caption{Ánh xạ 4 lớp IoT chuẩn sang thành phần thực trong repo.}\\
\toprule
\textbf{Lớp} & \textbf{Vai trò (ẩn dụ)} & \textbf{Thành phần phần cứng/mạng} & \textbf{Thư mục chính trong repo} \\
\midrule
\endhead
\textbf{Cảm biến} (Perception) & Giác quan: đo thực tế, tạo gói dữ liệu &
6\,$\times$ JSN-SR04T + ESP32-S3 sensor-node + còi GPIO11 &
\code{firmware/shared/thresholds.h}, \code{firmware/sensor-node/src/} \\
\textbf{Mạng} (Network) & Hệ thần kinh: vận chuyển dữ liệu &
ESP-NOW (đường chính), Wi-Fi \& MQTT (đường phụ) &
\code{espnow\_protocol.h}, \code{espnow\_client.cpp}, \code{espnow\_receiver.c}, \code{plugins/coreiot} \\
\textbf{Xử lý} (Processing) & Não bộ: lưu trữ, phân loại, ra quyết định &
CoreIoT broker chạy Rule-Chain (ThingsBoard) &
\code{cloud/coreiot/rule\_chain/}, \code{tools/test\_mqtt\_coreiot.py} \\
\textbf{Ứng dụng} (Application) & Ý thức: hiển thị, tự động hoá &
Waveshare 7" RGB LVGL v9 + Dashboard cloud &
\code{firmware/waveshare-screen/src/}, \code{components/ui\_dashboard}, \code{components/coreiot\_client} \\
\bottomrule
\end{longtable}
\end{center}

\section{Hai đường truyền song song (Hybrid V2)}

Điểm khác biệt so với mô hình IoT 4 lớp "dọc" thông thường: hệ thống này có
\textbf{hai tuyến} cùng vận chuyển gói dữ liệu \code{espnow\_sensor\_msg\_t}:

\begin{center}
\begin{tikzpicture}[
  node distance=1.4cm and 2.0cm,
  box/.style={draw, rounded corners=2pt, minimum height=0.9cm, minimum width=3.2cm, font=\small, align=center},
  lbl/.style={font=\small\bfseries, align=center}
]
  % LỚP 1
  \node[lbl] (l1) at (0,0) {Lớp 1 — Cảm biến};
  \node[box, fill=blue!8] (sen) at (0,-1.1) {6$\times$ JSN-SR04T\\ESP32-S3 sensor-node\\SensorTask (100 ms)};
  % LỚP 2
  \node[lbl] (l2) at (0,-2.8) {Lớp 2 — Mạng};
  \node[box, fill=green!8] (espnow) at (-3.2,-4.0) {ESP-NOW\\kênh 1, 500 ms\\đường chính <50 ms};
  \node[box, fill=green!12] (mqtt) at (3.2,-4.0) {MQTT/Wi-Fi\\app.coreiot.io:1883\\đường phụ 100--500 ms};
  % LỚP 3
  \node[lbl] (l3) at (3.2,-5.8) {Lớp 3 — Xử lý (CoreIoT/ThingsBoard)};
  \node[box, fill=orange!12] (rc) at (3.2,-7.0) {Rule-Chain 7 node\\script JS (100/30 cm)\\Timeseries + Attributes};
  % LỚP 4
  \node[lbl] (l4) at (0,-8.8) {Lớp 4 — Ứng dụng};
  \node[box, fill=red!8] (screen) at (-3.2,-10.0) {Waveshare-screen\\LVGL v9 UI\\Cảnh báo + còi hiển thị};
  \node[box, fill=red!12] (dash) at (3.2,-10.0) {Dashboard cloud\\warning\_status\\B9 nghiệm thu};
  % ARROWS
  \draw[-{Stealth}] (sen) -- (espnow) node[midway, right, font=\small\itshape] {packed 30 B};
  \draw[-{Stealth}] (sen) -- (mqtt) node[midway, right, font=\small\itshape] {telemetry 500 ms};
  \draw[-{Stealth}] (espnow) -- (screen) node[midway, right, font=\small\itshape] {trực tiếp};
  \draw[-{Stealth}] (mqtt) -- (rc);
  \draw[-{Stealth}] (rc) -- node[midway, right, align=center, font=\small\itshape] {attributes\\notifyDevice} (screen);
  \draw[-{Stealth}] (rc) -- (dash);
\end{tikzpicture}
\captionof{figure}{Hành trình của một gói dữ liệu khoảng cách qua 2 đường truyền.}
\end{center}

\textbf{Đường chính (ESP-NOW)}: SensorTask đọc/lọc 6 cảm biến $\rightarrow$
\code{espnow\_client} gửi gói 30\,B mỗi 500\,ms $\rightarrow$
\code{espnow\_receiver} trên screen nhận, cập nhật \code{sensor\_model}
$\rightarrow$ UI đổi màu. Toàn tuyến mục tiêu dưới 50\,ms và \textbf{không phụ
thuộc Internet}.

\textbf{Đường phụ (MQTT/CoreIoT)}: \code{coreiot\_client} publish telemetry
\texttt{d1..d6 + nearest\_cm} mỗi 500\,ms $\rightarrow$ broker
\code{app.coreiot.io} $\rightarrow$ Rule-Chain chạy script JS phân loại
\code{warning\_status} $\rightarrow$ lưu Timeseries, đẩy Attributes ngược về
screen và hiển thị Dashboard.

\section{Cây repo theo lớp}

\begin{verbatim}
UMT-TBS-official/
|-- firmware/
|   |-- shared/                 [HOP DONG (R2/R3/R4) - dung cho ca 2 board]
|   |   |-- espnow_protocol.h   [Lop 2: payload ESP-NOW packed]
|   |   +-- thresholds.h        [Lop 1: cham cam bien, nguong, buzzer]
|   |-- sensor-node/            [Lop 1 + 2: đo, loc, gui]
|   |   `-- src/
|   |       |-- ultrasonic_sensor.cpp   [1: đọc Trig/Echo]
|   |       |-- distance_filter.cpp     [1: loc cluster-EMA]
|   |       |-- shared_state.cpp        [1: trang thai 6 slot]
|   |       |-- buzzer.cpp              [1: coi GPIO11]
|   |       |-- espnow_client.cpp       [2: gui ESP-NOW]
|   |       |-- main.cpp                [1: SensorTask/BuzzerTask/NetworkTask]
|   |       `-- plugins/coreiot/
|   |           `-- coreiot_client.*    [2: MQTT 500 ms]
|   `-- waveshare-screen/       [Lop 4: hien thi, nhan 2 duong]
|       |-- src/main.c + bsp/           [4: khoi dong, RGB LCD + GT911]
|       `-- components/
|           |-- espnow_receiver/        [2: nhan ESP-NOW, LINKED badge]
|           |-- sensor_model/           [4: trang thai + classify + stale]
|           |-- ui_dashboard/           [4: giao dien LVGL v9]
|           `-- coreiot_client/         [2+4: MQTT nhan attributes]
|-- cloud/coreiot/rule_chain/   [Lop 3: snapshot Rule-Chain (R11)]
|-- tools/                      [Lop 3: test MQTT, guard scripts]
|-- .github/workflows/ci.yml    [CI + bao mat]
`-- report-code/                [ban bao cao nay]
\end{verbatim}

\section{Quy ước đọc báo cáo}

\begin{itemize}
\item \code{\textbackslash code\{...\}} — tên file/ký hiệu trong code.
\item Mỗi mục hàm có bảng: \textbf{Ký hiệu} / \textbf{File} / \textbf{Lớp} /
\textbf{Vai trò}, rồi phân tích dòng-cuối có tham chiếu dòng thực tế
(\code{file.cpp:34}).
\item Gắn cờ quy tắc CONSTITUTION khi liên quan: \textbf{R2} shared contract,
\textbf{R3} thresholds một nguồn, \textbf{R4} SENSOR\_COUNT từ
\code{sizeof}, \textbf{R7} file $\le$ 400 dòng, \textbf{R11} rule-chain
snapshot.
\item Ảnh raster nếu có đặt trong \code{report-code/pictures/} (macro
\code{\textbackslash pic}); sơ đồ vector vẽ bằng TikZ ngay trong chương.
\end{itemize}